\documentclass[12pt]{article}
\usepackage[utf8]{inputenc}
\usepackage{amsmath}
\usepackage{amsfonts}
\usepackage{amssymb}
\usepackage{geometry}
\usepackage{booktabs}
\usepackage{array}
\usepackage{float}

\geometry{margin=1in}

\title{Delay Costs Methodology\\Comprehensive Transmission Cost Calculator}
\author{Andrew Igdal}
\date{\today}

\begin{document}

\maketitle

\section{Overview}
This document describes the methodology for calculating delay costs in transmission line construction projects. Delay costs represent the additional expenses incurred when project timelines are extended beyond the original schedule.

\section{Methodology}

\subsection{Delay Cost Categories}
The delay cost calculation considers the following categories of costs that accumulate during project delays:

\begin{enumerate}
    \item \textbf{Legal Delay Costs}: Legal fees and litigation costs associated with delays
    \item \textbf{Administrative Overhead Delay Costs}: Administrative expenses during extended project timeline
    \item \textbf{Engineering Delay Costs}: Additional engineering work and redesign costs
    \item \textbf{Permitting and Regulatory Compliance Delay Costs}: Extended permitting and compliance expenses
    \item \textbf{Material and Equipment Delay Costs}: Storage, depreciation, and replacement costs for delayed materials
    \item \textbf{Labor Delay Costs}: Extended labor costs and workforce management expenses
    \item \textbf{Public Relations Delay Costs}: Communication and stakeholder management during delays
    \item \textbf{Environmental Delay Costs}: Additional environmental monitoring and mitigation costs
    \item \textbf{Miscellaneous Delay Costs}: Other unforeseen costs associated with project delays
\end{enumerate}

\subsection{Mathematical Formulation}

The total annual delay costs are calculated as:
\begin{equation}
C_{delay,total} = \sum_{i=1}^{9} C_{delay,i}
\end{equation}

Where $C_{delay,i}$ represents the annual cost for delay category $i$.

The present value of total delay costs over the delay timeline is calculated as:
\begin{equation}
PV_{delay} = \sum_{t=1}^{T} \frac{C_{delay,total}}{(1 + r)^t}
\end{equation}

Where:
\begin{itemize}
    \item $T$ = Delay timeline in years
    \item $r$ = Firm interest rate (decimal)
    \item $t$ = Time period (year)
\end{itemize}

\section{Input Parameters}

\begin{table}[H]
\centering
\caption{Required Input Parameters}
\begin{tabular}{@{}ll@{}}
\toprule
\textbf{Parameter} & \textbf{Description} \\
\midrule
Delay Timeline & Duration of project delay in years \\
Legal Delay Costs & Annual legal costs during delay (USD) \\
Administrative Overhead Delay Costs & Annual administrative costs during delay (USD) \\
Engineering Delay Costs & Annual engineering costs during delay (USD) \\
Permitting Delay Costs & Annual permitting costs during delay (USD) \\
Material and Equipment Delay Costs & Annual material/equipment costs during delay (USD) \\
Labor Delay Costs & Annual labor costs during delay (USD) \\
Public Relations Delay Costs & Annual PR costs during delay (USD) \\
Environmental Delay Costs & Annual environmental costs during delay (USD) \\
Miscellaneous Delay Costs & Annual miscellaneous costs during delay (USD) \\
Firm Interest Rate & Discount rate for present value calculation (decimal) \\
\bottomrule
\end{tabular}
\end{table}

\section{Outputs}

The script provides the following outputs:
\begin{itemize}
    \item Total annual delay costs
    \item Present value of total delay costs over the project timeline
\end{itemize}

\section{Assumptions}

\begin{itemize}
    \item Delay costs are incurred annually throughout the delay period
    \item The firm interest rate remains constant over the delay timeline
    \item All costs are expressed in nominal USD
    \item Delay costs are additive across all categories
\end{itemize}

\section{Usage Notes}

\begin{itemize}
    \item Input all cost values in USD
    \item Interest rate should be entered as a decimal (e.g., 0.05 for 5\%)
    \item Delay timeline should be entered in years
    \item All delay cost categories should be entered as annual costs
\end{itemize}

\end{document}
